\chapter{آليات تفاعلات المعقدات}\label{ch:b3:complex-mechanisms}

أذِب ملحًا للكروم(III) في ماء موسوم بالأكسجين-18، فتمضي أيام
قبل أن يحل الماء الموسوم محل جزيئات الماء المرتبطة
بالفلز؛ أما مع ملح للنحاس(II) فيكتمل التبادل في زمن
الخلط. والتفاعل هو نفسه، جزيء ماء يغادر أيونًا فلزيًا 
وآخر يأخذ مكانه، ومع ذلك تتغير سرعته تغيرًا هائلًا من أيون إلى
آخر، ويتنبأ عدد إلكترونات $d$ بجزء كبير من هذا التغير. وطريقة
استبدال الربائط هي التي تحدد كيف تُصنع أدوية البلاتين المضادة للسرطان وكيف
تعمل؛ وطريقة انتقال الإلكترونات بين الأيونات الفلزية هي التي تحدد سرعة
التنفس وسرعة كل بطارية. ويعالج هذا الفصل عائلتي
تفاعلات المعقدات: الاستبدال وانتقال الإلكترون، والثانية مع
نظرية رودولف ماركوس.

\begin{recall}
وصف مجلد السنة الثانية المعقدات، ودرجة تسنين الربائط، والربائط
الجسرية، والمتماكبات $fac$ و$mer$، وتبادل الربائط خطوةً أولية، 
وطاقات استقرار المجال البلوري مع التشكيلات عالية السبين ومنخفضة السبين؛
ووصف مجلد السنة الأولى الآليات مع تقريبي الحالة المستقرة والخطوة المحددة
للسرعة. وأعطى \cref{ch:b3:rate-theories} \omterm{thm:b3:rate-theories:eyring}{معادلة أيرنغ} 
ومعنى $\Delta^{\ddagger}S^\circ$، و\cref{ch:b3:electrode-kinetics} انتقال الإلكترون عند
القطب، و\cref{ch:b3:complex-spectra-magnetism} \omterm{def:b3:complex-spectra-magnetism:ligand-field-term}{حدود مجال الربائط} 
ومفعول يان--تيلر.
\end{recall}

\section{سرعة الاستبدال والخمول الحركي}

\begin{definition}[المعقدات سريعة الاستبدال والخاملة]\label{def:b3:complex-mechanisms:labile}
يكون المعقد \emph{سريع الاستبدال}\index{معقد سريع الاستبدال} إذا استُبدلت ربائطه
بسرعة، و\emph{معقدًا خاملًا حركيًا}\index{معقد خامل حركيًا}
إذا استُبدلت ببطء؛ وتبعًا لتاوبه، يسمى المعقد سريع الاستبدال حين
تنتهي تفاعلاته مع ربيطة تركيزها \qty{0.1}{mol/L} عند \qty{25}{\celsius} في نحو
دقيقة. وسرعة الاستبدال خاصية حركية، لا علاقة لها بالاستقرار الترموديناميكي
للمعقد.
\end{definition}

تبادل الماء، $\mathrm{[M(H_2O)_6]^{n+} + H_2O^{\ast} \to [M(H_2O)_5(H_2O^{\ast})]^{n+} + H_2O}$، أبسط
استبدال وهو يضع سلّمًا. فأيونات المجموعات الرئيسة تتبادل أسرع كلما
كانت أكبر وأقل شحنة. ومن بين أيونات الفلزات الانتقالية، تكون الأيونات ذات التشكيل $d^3$
(\ce{Cr^{3+}}) أو $d^6$ منخفض السبين (\ce{Co^{3+}}، \ce{Rh^{3+}}، \ce{Ir^{3+}}) خاملة، وكذلك
معظم فلزات الصفين الثاني والثالث، بينما تكون $d^9$ (\ce{Cu^{2+}}) و$d^4$ عالي السبين (\ce{Cr^{2+}})،
المشوَّهة بمفعول يان--تيلر، سريعة الاستبدال إلى أقصى حد.

\begin{proposition}[طاقة التنشيط لمجال الربائط]\label{prop:b3:complex-mechanisms:lfae}
إذا مر استبدال بحالة انتقالية خماسية التناسق على شكل هرم مربع القاعدة،
فإن فقد استقرار مجال الربائط عند بلوغها (طاقة التنشيط
لمجال الربائط) يكون، في أبسط نموذج $\sigma$ وحده، أكبر ما يكون من أجل $d^6$ منخفض السبين
($4\,Dq$) ثم من أجل $d^3$ و$d^8$ ($2\,Dq$)، ومعدومًا من أجل $d^0$ و$d^5$ عالي السبين و$d^{10}$، وسالبًا
من أجل $d^4$ و$d^9$ عاليي السبين.
\end{proposition}

\begin{proof}
لننمذج كل ربيطة على أنها ترفع مدارًا $d$ بمقدار $e_\sigma$ مضروبًا في مربع السعة
الزاوية للمدار في اتجاهها، مع تسوية القيمة إلى 1 على محور الفص: فالربيطة
المحورية تعطي $d_{z^2}$ القيمة 1، والاستوائية تعطي $d_{z^2}$ القيمة $\frac14$ و$d_{x^2-y^2}$ القيمة $\frac34$؛ والمدارات $t_{2g}$ تتجه
بين الربائط فلا تنال شيئًا. ويعطي ثماني الأوجه كلًّا من $d_{z^2}$ و$d_{x^2-y^2}$ القيمة $3e_\sigma$
($\Delta_{\mathrm o} = 3e_\sigma = 10Dq$)؛ ويعطي الهرم المربع، دون إحدى الربيطتين المحوريتين، $d_{z^2}$ القيمة $2e_\sigma$ 
و$d_{x^2-y^2}$ القيمة $3e_\sigma$. واستقرار $n$ إلكترونًا هو طاقتها مطروحًا منها $n/5$ مرة مجموع
المستويات الخمسة (مركز الثقل). ومن أجل $d^6$ منخفض السبين ($t_{2g}^6$): $0 - \frac65 \times 6e_\sigma$ في
ثماني الأوجه، و$0 - \frac65 \times 5e_\sigma$ في الهرم: أي فقد قدره $1.2e_\sigma = 4Dq$. ومن أجل $d^3$: $\frac35e_\sigma =
2Dq$؛ ومن أجل $d^8$ ($t_{2g}^6e_g^2$): $(5 - \frac85 \times 5) - (6 - \frac85 \times 6) = 0.6e_\sigma$. ومن أجل $d^9$: $(7 - 9) - (9 -
10.8) = -0.2e_\sigma$.
\end{proof}

\begin{omfigure}
\begin{tikzpicture}
\begin{axis}[width=0.72\linewidth, height=5.4cm, xmin=-0.6, xmax=10.6, ymin=-1.2, ymax=4.6,
  axis x line=bottom, axis y line=left, ycomb, xtick={0,1,2,3,4,5,6,7,8,9,10},
  xlabel={\babelsublr{$n$ ($d^n$)}}, ylabel={\foreignlanguage{arabic}{طاقة التنشيط بوحدة $Dq$}}, tick label style={font=\small},
  label style={font=\small}]
\draw[black!40] (axis cs:-0.6,0) -- (axis cs:10.6,0);
\addplot[omDef, line width=4pt, xshift=-2.5pt] table[x=n, y=hs] {figdata/out/bachelor-3-ligand-field-activation.dat};
\addplot[omThm, line width=4pt, xshift=2.5pt, unbounded coords=discard] table[x=n, y=ls] {figdata/out/bachelor-3-ligand-field-activation.dat};
\end{axis}
\end{tikzpicture}

{\small طاقات التنشيط لمجال الربائط من أجل طريق تفككي عبر
هرم مربع، في نموذج $\sigma$ وحده الوارد في القضية (بالأزرق: عالي السبين؛
بالأحمر: منخفض السبين). والقيم الأكبر،
من أجل $d^6$ منخفض السبين و$d^3$ و$d^8$، توافق الأيونات الخاملة أو البطيئة (\ce{Co^{3+}}، \ce{Cr^{3+}}، \ce{Ni^{2+}})؛
والقيم السالبة، من أجل $d^4$ و$d^9$، توافق أيونات يان--تيلر السريعة الاستبدال جدًا.}
\end{omfigure}

\begin{definition}[حجم التنشيط]\label{def:b3:complex-mechanisms:activation-volume}
\emph{حجم التنشيط}\index{حجم التنشيط} $\Delta V^{\ddagger}$ لخطوة أولية
هو الفرق بين الحجم المولي الجزئي \omterm{def:b3:rate-theories:activated-complex}{للمعقد
المنشَّط} والحجوم المولية الجزئية للمتفاعلات.
\end{definition}

\begin{proposition}[تأثير الضغط في ثابت السرعة]\label{prop:b3:complex-mechanisms:pressure}
عند درجة حرارة ثابتة، $\displaystyle\Bigl(\frac{\partial\ln k}{\partial p}\Bigr)_T = -\frac{\Delta V^{\ddagger}}{RT}$.
\end{proposition}

\begin{proof}
بحسب \omterm{thm:b3:rate-theories:eyring}{معادلة أيرنغ} $\ln k = \text{ثابت} - \Delta^{\ddagger}G/RT$، و$(\partial\Delta^{\ddagger}G/\partial p)_T = \Delta V^{\ddagger}$، كما هو الشأن في
كل طاقة جيبس.
\end{proof}

للخطوة التي تغادر فيها ربيطة $\Delta V^{\ddagger} > 0$؛ وللخطوة التي تدخل فيها ربيطة $\Delta V^{\ddagger}
< 0$. \omterm{def:b3:complex-mechanisms:activation-volume}{وحجم التنشيط}، المقيس عند بضع مئات من الميغاباسكال، هو أكثر
وسائل تشخيص آلية الاستبدال مباشرة، وهو أقل تشويشًا بالإذابة من
\omterm{def:b3:rate-theories:activation}{إنتروبيا التنشيط}.

\section{آليات الاستبدال}

\begin{definition}[الآليات التفككية والترابطية والتبادلية]\label{def:b3:complex-mechanisms:mechanisms}
في \emph{الآلية التفككية}\index{آلية تفككية} (D) تغادر الربيطة
المغادرة أولًا، فتعطي وسيطًا عدد تناسقه أدنى؛ وفي
\emph{الآلية الترابطية}\index{آلية ترابطية} (A) ترتبط الربيطة الداخلة
أولًا، فتعطي وسيطًا عدد تناسقه أعلى؛ وفي
\emph{الآلية التبادلية}\index{آلية تبادلية} (I) تتبادل الربيطتان
مكانيهما في خطوة واحدة، بطابع يغلب عليه كسر الرابطة ($I_d$) أو تكوينها ($I_a$).
\end{definition}

\begin{proposition}[قانونا السرعة للآليتين D وA]\label{prop:b3:complex-mechanisms:d-a-laws}
من أجل $\mathrm{ML_5X + Y \to ML_5Y + X}$: تعطي الآلية D ($\mathrm{ML_5X} \rightleftharpoons \mathrm{ML_5} + \mathrm X$، $k_1$، $k_{-1}$؛ $\mathrm{ML_5} + \mathrm Y
\to \mathrm{ML_5Y}$، $k_2$)
\[
  v = \frac{k_1k_2[\mathrm{ML_5X}][\mathrm Y]}{k_{-1}[\mathrm X] + k_2[\mathrm Y]},
\]
وهو من الرتبة الأولى بالنسبة إلى المعقد ومستقل عن $[\mathrm Y]$ عند $[\mathrm Y]$ مرتفع؛ أما الآلية A
عبر وسيط سباعي التناسق في حالة مستقرة فتعطي $v =
k[\mathrm{ML_5X}][\mathrm Y]$، من الرتبة الأولى بالنسبة إلى كل منهما.
\end{proposition}

\begin{proof}
D: الحالة المستقرة من أجل $\mathrm{ML_5}$، $k_1[\mathrm{ML_5X}] = (k_{-1}[\mathrm X] + k_2[\mathrm Y])[\mathrm{ML_5}]$، و$v = k_2[\mathrm{ML_5}][\mathrm Y]$.
وعند $[\mathrm Y]$ مرتفع، $v \to k_1[\mathrm{ML_5X}]$. A: $\mathrm{ML_5X} + \mathrm Y \rightleftharpoons \mathrm{ML_5XY}$ ($k_a$، $k_{-a}$)، $\mathrm{ML_5XY} \to
\mathrm{ML_5Y} + \mathrm X$ ($k_b$)؛ والحالة المستقرة تعطي $v = \frac{k_ak_b}{k_{-a} + k_b}[\mathrm{ML_5X}][\mathrm Y]$.
\end{proof}

في الماء، تكون الربيطة الداخلة عادة بتركيز أدنى بكثير من تركيز
المذيب، والخطوة الأولى الحقيقية هي التقاء المعقد 
والربيطة.

\begin{definition}[آلية آيغن--ويلكنز]\label{def:b3:complex-mechanisms:eigen-wilkins}
\emph{آلية آيغن--ويلكنز}\index{آلية آيغن--ويلكنز} لاستبدال
معقد مائي هي توازن مسبق سريع يكوّن
\emph{معقد الكرة الخارجية}\index{معقد الكرة الخارجية}، الذي تجلس فيه الربيطة
الداخلة بجوار المعقد، خارج كرة تناسقه، يليه
تبادل جزيء ماء بتلك الربيطة، وهو الخطوة المحددة للسرعة.
\end{definition}

\begin{theorem}[قانون سرعة آيغن--ويلكنز]\label{thm:b3:complex-mechanisms:eigen-wilkins}
مع ثابت توازن الكرة الخارجية $K_{\mathrm{os}}$ وثابت سرعة التبادل
$k_i$، يكون ثابت السرعة الملاحظ من الرتبة الأولى الزائفة عند فائض من $[\mathrm Y]$
\[
  k_{\mathrm{obs}} = \frac{K_{\mathrm{os}}k_i[\mathrm Y]}{1 + K_{\mathrm{os}}[\mathrm Y]} .
\]
\end{theorem}

\begin{proof}
يتوزع المعقد بين الصورة الحرة وصورة الكرة الخارجية، $[\mathrm{M\cdot Y}] = K_{\mathrm{os}}[\mathrm M][\mathrm Y]$،
والمجموع $[\mathrm M]_{\mathrm{tot}} = [\mathrm M](1 + K_{\mathrm{os}}[\mathrm Y])$؛ والسرعة $k_i[\mathrm{M\cdot Y}] = k_iK_{\mathrm{os}}[\mathrm Y][\mathrm M]_{\mathrm{tot}}/(1 + K_{\mathrm{os}}[\mathrm Y])$.
\end{proof}

عند $[\mathrm Y]$ منخفض يكون التفاعل من الرتبة الثانية مع $k = K_{\mathrm{os}}k_i$؛ و$K_{\mathrm{os}}$، المقدَّر من
الشحنات ومن مسافة الاقتراب الأقصى، يكاد يكون نفسه لكل
الربائط ذات الشحنة نفسها، بحيث إن $k_i$، أي خطوة تبادل الماء، هو الذي يتحكم في
السرعة: فالاستبدال على أيون مائي يكاد يكون مستقلًا عن الربيطة
الداخلة، وهذه بصمة التبادل التفككي.

\begin{definition}[آلية القاعدة المرافقة]\label{def:b3:complex-mechanisms:conjugate-base}
\emph{آلية القاعدة المرافقة}\index{آلية القاعدة المرافقة} للحلمأة
القاعدية لمعقد أميني، $\mathrm{[Co(NH_3)_5X]^{2+} + OH^- \to [Co(NH_3)_5(OH)]^{2+} + X^-}$، هي
نزع بروتون من ربيطة أمين في توازن مسبق سريع، يليه
فقد تفككي للأيون $\mathrm X^-$ من القاعدة المرافقة الأميدية، التي تجعل ربيطتها $\mathrm{NH_2^-}$
المعقدَ \omterm{def:b3:complex-mechanisms:labile}{سريع الاستبدال}.
\end{definition}

\begin{proposition}[قانون سرعة الحلمأة القاعدية]\label{prop:b3:complex-mechanisms:conjugate-base}
مع ثابت نزع البروتون $K$ وثابت سرعة التفكك $k$ للقاعدة
المرافقة، $v = \dfrac{kK[\text{المعقد}]_{\mathrm{tot}}[\ce{OH-}]}{1 + K[\ce{OH-}]}$، من الرتبة الثانية عند $[\ce{OH-}]$ منخفض.
\end{proposition}

\begin{proof}
كما في قانون آيغن--ويلكنز: يتوزع المعقد بين صورته الحمضية 
وصورة قاعدته المرافقة بالنسبة $1 : K[\ce{OH-}]$، والثانية وحدها تتفاعل.
\end{proof}

لا يميّز قانون السرعة وحده هذه الآلية من هجوم مباشر للأيون
\ce{OH-}؛ والدليل أن الحلمأة القاعدية تحتاج إلى بروتون N--H (فالمعقدات التي تخلو
منه لا تُظهرها) وأن بروتونات الأمين تتبادل مع المذيب أسرع بكثير
من الحلمأة.

\begin{method}[تشخيص آلية استبدال]\label{met:b3:complex-mechanisms:diagnose}
\begin{enumerate}
  \item الاعتماد على المجموعة الداخلة: سرعة مستقلة عن Y (أو بلوغ إشباع)،
        \omterm{def:b3:complex-mechanisms:mechanisms}{فالآلية تفككية}؛ وسرعة متناسبة مع $[\mathrm Y]$ وحساسة لطبيعة Y،
        \omterm{def:b3:complex-mechanisms:mechanisms}{فالآلية ترابطية}.
  \item \omterm{def:b3:complex-mechanisms:activation-volume}{حجم التنشيط}: موجب، فهي تفككية؛ سالب، فهي ترابطية.
  \item \omterm{def:b3:rate-theories:activation}{إنتروبيا التنشيط}: موجبة في D، وسالبة في A (بحذر، لأن
        الإذابة تُسهم فيها).
  \item الاعتماد على المجموعة المغادرة: سرعة تتبع قوة الرابطة M--X
        تشير إلى كسر الرابطة في الحالة الانتقالية.
\end{enumerate}
\end{method}

\begin{omfigure}
\begin{tikzpicture}
\begin{axis}[name=pd, omprofile, width=0.46\linewidth, height=4.6cm, xmin=0, xmax=10, ymin=0, ymax=10,
  xlabel={\foreignlanguage{arabic}{إحداثي التفاعل}}, ylabel={\foreignlanguage{arabic}{الطاقة}}, title={\small \babelsublr{D: $\mathrm{ML_5X} \to \mathrm{ML_5} + \mathrm X \to \mathrm{ML_5Y}$}}]
\addplot[omDef, very thick, smooth] coordinates {(0,2) (1.5,2.1) (3,7.5) (4.2,5.6) (5.0,5.4) (5.8,5.6) (7,7.0) (8.5,1.6) (10,1.5)};
\node[font=\scriptsize, anchor=north, align=center] at (axis cs:5,5.2) {\foreignlanguage{arabic}{\babelsublr{ML$_5$}، تناسق 5}};
\end{axis}
\begin{axis}[at={(pd.east)}, anchor=west, xshift=1.0cm, omprofile, width=0.46\linewidth, height=4.6cm,
  xmin=0, xmax=10, ymin=0, ymax=10, xlabel={\foreignlanguage{arabic}{إحداثي التفاعل}}, ylabel={\foreignlanguage{arabic}{الطاقة}},
  title={\small \babelsublr{A: $\mathrm{ML_5X} + \mathrm Y \to \mathrm{ML_5XY} \to \mathrm{ML_5Y}$}}]
\addplot[omThm, very thick, smooth] coordinates {(0,2) (1.5,2.1) (3,6.5) (4.2,4.4) (5.0,4.2) (5.8,4.4) (7,7.2) (8.5,1.6) (10,1.5)};
\node[font=\scriptsize, anchor=north, align=center] at (axis cs:5,4.0) {\foreignlanguage{arabic}{\babelsublr{ML$_5$XY}، تناسق 7}};
\end{axis}
\end{tikzpicture}

{\small مخططا الطاقة لاستبدال تفككي واستبدال ترابطي
(رسم تخطيطي): كلاهما يمر بوسيط، عدد تناسقه أدنى
في D وأعلى في A. وفي الآليات التبادلية
تختفي البئر ولا تبقى إلا حالة انتقالية واحدة.}
\end{omfigure}

المعقد الثماني الأوجه الذي يُستبدل عبر هرم مربع يحافظ عمومًا
على تشكيله ($cis$ يبقى $cis$)؛ أما الذي يمر عبر هرم ثنائي
مثلثي، كما تفعل معقدات كثيرة للكوبالت(III) في الحلمأة القاعدية، فقد يعطي
مزيجًا من نواتج $cis$ و$trans$.

\section{الاستبدال في المعقدات المربعة المستوية ومفعول ترانس}

المعقدات المربعة المستوية لأيونات $d^8$ (\ce{Pt^{2+}}، \ce{Pd^{2+}}، \ce{Au^{3+}}، \ce{Ni^{2+}} مع ربائط قوية)
غير مشبعة تناسقيًا: إذ تستطيع ربيطة داخلة أن تقترب على امتداد
الاتجاه المحوري الفارغ. واستبدالاتها ترابطية، عبر حالة انتقالية
خماسية التناسق على شكل هرم ثنائي مثلثي.

\begin{proposition}[قانون السرعة ذو الحدين]\label{prop:b3:complex-mechanisms:square-planar}
يخضع الاستبدال $\mathrm{[PtL_3X] + Y \to [PtL_3Y] + X}$ في مذيب متناسق S
للعلاقة $k_{\mathrm{obs}} = k_1 + k_2[\mathrm Y]$ عند فائض من Y.
\end{proposition}

\begin{proof}
طريقان ترابطيان متوازيان: هجوم مباشر للربيطة Y ($k_2[\mathrm Y]$)، وهجوم
المذيب، الموجود بفائض كبير (من الرتبة الأولى الزائفة $k_1$)، فيعطي $[\mathrm{PtL_3S}]$، الذي
يحوّله Y بعد ذلك بسرعة. والطرق المتوازية من الرتبة الأولى تُجمع.
\end{proof}

\begin{definition}[مفعول ترانس، تأثير ترانس]\label{def:b3:complex-mechanisms:trans-effect}
\emph{مفعول ترانس}\index{مفعول ترانس} هو تسريع استبدال
ربيطة بفعل الربيطة الواقعة في موضع ترانس منها: وهو مفعول حركي، يخص الحالة الانتقالية.
و\emph{تأثير ترانس}\index{تأثير ترانس} هو إضعاف الرابطة الواقعة في موضع ترانس من ربيطة، في الحالة
الأساسية، ويُرى في أطوال الروابط وفي الترددات
الاهتزازية: وهو مفعول ترموديناميكي.
\end{definition}

يتبع \omterm{def:b3:complex-mechanisms:trans-effect}{مفعول ترانس} الترتيب $\ce{CO}$، $\ce{CN-}$، $\ce{C2H4} > \ce{PR3}$، $\ce{H-} > \ce{CH3-} > \ce{I-}$، $\ce{SCN-} >
\ce{Br-} > \ce{Cl-} > \ce{NH3}$، $\ce{H2O}$. فالمانحات $\sigma$ القوية تُضعف الرابطة المقابلة (\omterm{def:b3:complex-mechanisms:trans-effect}{تأثير ترانس})؛
والمستقبلات $\pi$ القوية تثبّت الحالة الانتقالية الخماسية التناسق بأخذ
\omterm{def:b3:computational-chemistry:dft}{كثافة إلكترونية} من الفلز. وكلاهما يرفع السرعة.

\begin{method}[تخطيط تخليق متماكبات البلاتين(II)]\label{met:b3:complex-mechanisms:pt-synthesis}
\begin{enumerate}
  \item في كل خطوة، تكون الربيطة المستبدَلة هي الواقعة في موضع ترانس من الربيطة
        ذات \omterm{def:b3:complex-mechanisms:trans-effect}{مفعول ترانس} الأعلى بين الربائط الموجودة.
  \item رتّب الإضافات بحيث تقود هذه القاعدة إلى المتماكب المطلوب.
\end{enumerate}
\end{method}

\begin{omfigure}
\setchemfig{atom sep=2.2em}
\schemestart
\chemfig{Cl-Pt(-[2]Cl)(-[6]Cl)-Cl}
\arrow{->[\ce{NH3}]}
\chemfig{Cl-Pt(-[2]Cl)(-[6]Cl)-NH_3}
\arrow{->[\ce{NH3}]}
\chemfig{Cl-Pt(-[2]NH_3)(-[6]Cl)-NH_3}
\schemestop

\bigskip
\schemestart
\chemfig{H_3N-Pt(-[2]NH_3)(-[6]NH_3)-NH_3}
\arrow{->[\ce{Cl-}]}
\chemfig{H_3N-Pt(-[2]NH_3)(-[6]NH_3)-Cl}
\arrow{->[\ce{Cl-}]}
\chemfig{Cl-Pt(-[2]NH_3)(-[6]NH_3)-Cl}
\schemestop

{\small \omterm{def:b3:complex-mechanisms:trans-effect}{مفعول ترانس} في العمل (الشحنات محذوفة). في الأعلى: انطلاقًا من $\ce{[PtCl4]^{2-}}$، يحل
جزيء الأمونيا الثاني محل كلوريد مقابل لكلوريد (فلمفعول ترانس للأيون \babelsublr{Cl$^-$} قيمة أكبر)،
فيعطي المتماكب $cis$، أي السيسبلاتين. في الأسفل: انطلاقًا من $\ce{[Pt(NH3)4]^{2+}}$، يحل الكلوريد الثاني
محل الأمونيا المقابلة للكلوريد الأول، فيعطي المتماكب $trans$.}
\end{omfigure}

\section{انتقال الإلكترون}

\begin{definition}[انتقال الإلكترون عبر الكرة الخارجية وعبر الكرة الداخلية]\label{def:b3:complex-mechanisms:electron-transfer}
في \emph{انتقال الإلكترون عبر الكرة الخارجية}\index{انتقال الإلكترون عبر الكرة الخارجية} يمر
الإلكترون بين معقدين تبقى كرتا تناسقهما سليمتين؛
وفي \emph{انتقال الإلكترون عبر الكرة الداخلية}\index{انتقال الإلكترون عبر الكرة الداخلية}
يمر عبر ربيطة تجسر بين الفلزين في معقد ثنائي النواة
عابر. و\emph{تفاعل التبادل الذاتي}\index{تفاعل التبادل الذاتي} هو
انتقال إلكترون بين حالتي الأكسدة للمزدوجة نفسها، دون أي
تغير كيميائي صافٍ، مثل $\ce{[Fe(H2O)6]^{3+} + [Fe(H2O)6]^{2+}}$.
\end{definition}

بيّن هنري تاوبه الطريق عبر الكرة الداخلية في سنة 1953 بتفاعل اختير بحيث
يحكي الناتج قصته. فالكوبالت(III) خامل، والكوبالت(II) \omterm{def:b3:complex-mechanisms:labile}{سريع الاستبدال}؛ والكروم(II)
\omterm{def:b3:complex-mechanisms:labile}{سريع الاستبدال}، والكروم(III) خامل. وحين يُختزل $\ce{[Co(NH3)5Cl]^{2+}}$ بالمعقد $\ce{[Cr(H2O)6]^{2+}}$ في
وسط حمضي، يحمل كل الكروم(III) المتكوّن الكلوريد، في صورة $\ce{[Cr(H2O)5Cl]^{2+}}$، حتى في
وجود كلوريد حر: فلا بد أن الكلوريد جسَر بين الفلزين،
وارتبط بالكروم(III) الخامل فور مرور الإلكترون، 
وحرّره الكوبالت(II) \omterm{def:b3:complex-mechanisms:labile}{السريع الاستبدال}.

\begin{omfigure}
\begin{tikzpicture}[font=\small]
  \begin{scope}[scale=0.8]
    \pic[transform shape] (a) at (0,0) {omoctahedron};
    \pic[transform shape] (b) at (2.0,0) {omoctahedron};
  \end{scope}
  \node[anchor=north east] at (0.3,-1.0) {$\mathrm{Co^{III}(NH_3)_5}$};
  \node[anchor=north west] at (1.3,-1.0) {$\mathrm{Cr^{II}(H_2O)_5}$};
  \node[omThm, anchor=south] at (0.8,0.15) {Cl};
  \draw[->, thick, omDef] (1.6,1.75) to[bend right=40] node[midway, above] {e$^-$} (0,1.75);
\end{tikzpicture}

{\small الوسيط الجسري لتاوبه (رسم تخطيطي): الكلوريد، الذي لا يزال مرتبطًا
بالكوبالت(III)، يدخل كرة تناسق الكروم(II) بأن يحل محل جزيء
ماء؛ ويعبر الإلكترون عبر الجسر، ويبقى الكلوريد على
الكروم(III) الذي صار خاملًا.}
\end{omfigure}

يقفز الإلكترون في نحو $10^{-16}$ s، أسرع بكثير من حركة النوى: فبحسب
\omterm{thm:b3:electronic-spectroscopy:franck-condon}{مبدأ فرانك--كوندون} (\cref{ch:b3:electronic-spectroscopy})، لا يمكن أن يحدث الانتقال
إلا عند تشكيل نووي تكون فيه للمتفاعلات والنواتج
الطاقة نفسها. ففي التبادل الذاتي للحديد، الروابط \ce{Fe^{III}}--O أقصر من
الروابط \ce{Fe^{II}}--O؛ وقبل أن يستطيع الإلكترون الانتقال، يجب أن يتشوه المعقدان إلى
هندسة وسطى مشتركة، بكلفة طاقية.

\section{نظرية ماركوس}

\begin{definition}[طاقة إعادة التنظيم]\label{def:b3:complex-mechanisms:reorganisation}
\emph{طاقة إعادة التنظيم}\index{طاقة إعادة التنظيم} $\lambda$ لانتقال
إلكترون هي الطاقة اللازمة لجلب المتفاعلات، دون
نقل الإلكترون، إلى التشكيل النووي (أطوال روابط
المعقدات واتجاهات المذيب المحيط) الخاص بالنواتج. وهي
مجموع جزء داخلي، من الروابط، وجزء خارجي، من المذيب.
\end{definition}

\begin{theorem}[معادلة ماركوس]\label{thm:b3:complex-mechanisms:marcus}
إذا كانت طاقات جيبس للمتفاعلات والنواتج قطوعًا مكافئة متساوية
الانحناء في إحداثي تفاعل $x$ (0 عند توازن المتفاعلات، و1 عند توازن
النواتج)، فإن \omterm{def:b3:rate-theories:activation}{طاقة جيبس للتنشيط} لانتقال إلكترون
طاقة جيبس القياسية لتفاعله $\Delta G^\circ$ وطاقة إعادة تنظيمه $\lambda$ هي
\[
  \Delta G^{\ddagger} = \frac{(\lambda + \Delta G^\circ)^2}{4\lambda} .
\]
وهذه هي \emph{معادلة ماركوس}\index{معادلة ماركوس}.
\end{theorem}

\begin{proof}
اكتب $G_R = \lambda x^2$ و$G_P = \lambda(x - 1)^2 + \Delta G^\circ$: فالانحناء محدد بالعلاقة $G_R(1) = \lambda$، وهي
تعريف $\lambda$. ويحدث الانتقال حيث يتقاطع المنحنيان: $\lambda x^2 = \lambda x^2 - 2\lambda x + \lambda +
\Delta G^\circ$، ومنه $x^\ast = (\lambda + \Delta G^\circ)/2\lambda$، و$\Delta G^{\ddagger} = G_R(x^\ast) = (\lambda + \Delta G^\circ)^2/4\lambda$.
\end{proof}

\begin{corollary}[المنطقة المقلوبة]\label{cor:b3:complex-mechanisms:inverted}
عند $\lambda$ ثابتة، تتزايد السرعة كلما صار التفاعل أكثر إطلاقًا لطاقة جيبس حتى
$-\Delta G^\circ = \lambda$، حيث $\Delta G^{\ddagger} = 0$؛ وأبعد من ذلك، تتناقص من جديد.
\end{corollary}

\begin{proof}
$\Delta G^{\ddagger}$ قطع مكافئ في $\Delta G^\circ$ قيمته الدنيا، الصفر، عند $\Delta G^\circ = -\lambda$؛ ومن أجل $\Delta G^\circ < -\lambda$ يكبر
من جديد.
\end{proof}

\begin{definition}[المنطقة المقلوبة]\label{def:b3:complex-mechanisms:inverted}
\emph{المنطقة المقلوبة}\index{منطقة مقلوبة} لانتقال الإلكترون هي
مجال القوى الدافعة $-\Delta G^\circ > \lambda$ الذي تتناقص فيه السرعة كلما صار التفاعل
أكثر إطلاقًا لطاقة جيبس.
\end{definition}

\begin{omfigure}
\begin{tikzpicture}
\begin{axis}[name=mp, width=0.47\linewidth, height=5.6cm, xmin=-0.6, xmax=2.2, ymin=-160, ymax=120,
  axis lines=left, xlabel={\foreignlanguage{arabic}{إحداثي التفاعل $x$}}, ylabel={\babelsublr{$G$ / \unit{kJ.mol^{-1}}}},
  tick label style={font=\small}, label style={font=\small}, xtick={0,1,2}]
\addplot[black, thick] table[x=x, y=GR] {figdata/out/bachelor-3-marcus-parabolas.dat};
\addplot[omDef!50, thick] table[x=x, y=P0] {figdata/out/bachelor-3-marcus-parabolas.dat};
\addplot[omDef, thick] table[x=x, y=P50] {figdata/out/bachelor-3-marcus-parabolas.dat};
\addplot[omThm, thick] table[x=x, y=P100] {figdata/out/bachelor-3-marcus-parabolas.dat};
\addplot[omThm!60, thick, dashed] table[x=x, y=P150] {figdata/out/bachelor-3-marcus-parabolas.dat};
\node[font=\scriptsize, anchor=north] at (axis cs:0,-2) {R};
\end{axis}
\begin{axis}[at={(mp.east)}, anchor=west, xshift=1.4cm, width=0.47\linewidth, height=5.6cm,
  xmin=0, xmax=250, ymin=0, ymax=11.5, axis lines=left, xlabel={\babelsublr{$-\Delta G^\circ$ / \unit{kJ.mol^{-1}}}},
  ylabel={$\log(k/\unit{s^{-1}})$}, tick label style={font=\small}, label style={font=\small}]
\addplot[omDef, very thick] table[x=mdG, y=logk] {figdata/out/bachelor-3-marcus-rate.dat};
\draw[black!40, dashed] (axis cs:100,0) -- (axis cs:100,11);
\node[font=\scriptsize, anchor=south east] at (axis cs:95,1) {\foreignlanguage{arabic}{عادية}};
\node[font=\scriptsize, anchor=south west] at (axis cs:105,1) {\foreignlanguage{arabic}{مقلوبة}};
\end{axis}
\end{tikzpicture}

{\small نظرية ماركوس مع $\lambda = \qty{100}{kJ/mol}$ (نموذج). يسارًا: قطع المتفاعلات المكافئ
(بالأسود) وقطوع النواتج من أجل $\Delta G^\circ = 0$ و$-50$ (بالأزرق)، و$-100 = -\lambda$ (بالأحمر) 
و\qty{-150}{kJ/mol} (بالمتقطع، \omterm{def:b3:complex-mechanisms:inverted}{المنطقة المقلوبة})؛ ونقطة
تقاطعها، أي الحالة الانتقالية، تهبط إلى قاع منحنى المتفاعلات عند
$-\Delta G^\circ = \lambda$ وترتفع من جديد بعد ذلك. يمينًا: لوغاريتم ثابت السرعة
(بعامل قبلي نموذجي \qty{e11}{s^{-1}}) يمر بقيمة عظمى عند $-\Delta G^\circ = \lambda$.}
\end{omfigure}

\begin{proposition}[طاقة إعادة التنظيم الخارجية]\label{prop:b3:complex-mechanisms:outer-reorganisation}
من أجل متفاعلين كرويين نصفا قطريهما $a_1$ و$a_2$ على مسافة $d$ في مذيب
قرينة انكساره $n$ وسماحيته النسبية $\varepsilon_s$، يكون جزء المذيب من $\lambda$
متناسبًا مع $\bigl(\frac{1}{2a_1} + \frac{1}{2a_2} - \frac1d\bigr)\bigl(\frac{1}{n^2} - \frac{1}{\varepsilon_s}\bigr)$.
\end{proposition}

\admitted

يعطي الماء، الشديد القطبية ($\varepsilon_s$ كبيرة) لكن بقرينة انكسار عادية، إعادة
تنظيم كبيرة للمذيب؛ وتعطي المذيبات غير القطبية إعادة تنظيم صغيرة. والمتفاعلات الكبيرة
يُعاد تنظيمها أقل: ولهذا تدفن بروتينات نقل الإلكترون مواقعها الفلزية
داخل البروتين.

\begin{theorem}[علاقة ماركوس التقاطعية]\label{thm:b3:complex-mechanisms:cross-relation}
من أجل تفاعل تقاطعي $\mathrm{A_{ox} + B_{red} \to A_{red} + B_{ox}}$ ثابت توازنه $K_{12}$، بين مزدوجتين
ثابتا سرعة تبادلهما الذاتي $k_{11}$ و$k_{22}$،
\[
  k_{12} = \sqrt{k_{11}k_{22}K_{12}f}, \qquad \ln f = \frac{(\ln K_{12})^2}{4\ln(k_{11}k_{22}/Z^2)},
\]
مع $Z$ عامل تواتر التصادم؛ و$f \approx 1$ حين لا يكون $K_{12}$ كبيرًا جدًا. وهذه هي
\emph{علاقة ماركوس التقاطعية}\index{علاقة ماركوس التقاطعية}.
\end{theorem}

\begin{proof}[برهان جزئي]
لنفترض أن طاقة إعادة تنظيم التفاعل التقاطعي هي متوسط طاقتي
التبادلين الذاتيين، $\lambda_{12} = (\lambda_{11} + \lambda_{22})/2$ (فكل متفاعل يسهم بنصف \omterm{def:b3:complex-mechanisms:reorganisation}{طاقة إعادة التنظيم} في
تبادله الذاتي)، وأن كل ثوابت السرعة $k = Z\eu^{-\Delta G^{\ddagger}/RT}$. عندئذ
$\Delta G_{12}^{\ddagger} = \frac{\lambda_{12}}{4}\bigl(1 + \frac{\Delta G^\circ_{12}}{\lambda_{12}}\bigr)^2 = \frac{\lambda_{12}}{4} + \frac{\Delta G^\circ_{12}}{2} + \frac{(\Delta G^\circ_{12})^2}{4\lambda_{12}}$. ويعطي الحد الأول
$\sqrt{k_{11}k_{22}}$ (لأن $\Delta G^{\ddagger}_{ii} = \lambda_{ii}/4$)، والثاني $\sqrt{K_{12}}$ (مع $\Delta G^\circ_{12} = -RT\ln K_{12}$)، والثالث
العامل $f$.
\end{proof}

\begin{method}[سرعة تفاعل تقاطعي من معطيات التبادل الذاتي]\label{met:b3:complex-mechanisms:marcus-estimate}
\begin{enumerate}
  \item أوجد ثابتي التبادل الذاتي $k_{11}$ و$k_{22}$ للمزدوجتين.
  \item احسب $K_{12}$ من الجهود المعيارية: $\ln K_{12} = nF(E^\circ_{\text{المؤكسد}} - E^\circ_{\text{المختزل}})/RT$.
  \item $k_{12} = \sqrt{k_{11}k_{22}K_{12}}$ (مع $f$ إذا كان $K_{12}$ كبيرًا جدًا). والتوافق مع القيمة المقيسة في حدود عامل
        يساوي بضع مرات يدل على آلية عبر الكرة الخارجية.
\end{enumerate}
\end{method}

نشر ماركوس النظرية في سنة 1956؛ ورُصدت \omterm{def:b3:complex-mechanisms:inverted}{المنطقة المقلوبة}، المشكوك فيها أول الأمر،
في سنة 1984 على يد غيرهارد كلوس وجون ميلر في جزيئات يُثبَّت فيها المانح 
والمستقبل على مسافة ثابتة بفاصل صلب، بحيث لا يستطيع الانتشار
أن يخفيها.

\begin{inthelab}[تبادل الماء بالرنين المغناطيسي النووي للأكسجين-17]
يتعرض رنين الأكسجين-17 للماء المرتبط بأيون بارامغناطيسي ورنين
الماء الحر للتعريض بفعل التبادل بينهما. وتسجيل عرض خط
إشارة الماء الحر بدلالة درجة الحرارة، في محلول لبيركلورات الفلز
مُغنى بالنظير \ce{^{17}O}، يعطي ثابت سرعة التبادل، وفي مسبار عالي الضغط،
\omterm{def:b3:complex-mechanisms:activation-volume}{حجم تنشيطه}.
\end{inthelab}

\begin{safety}{\ghs[0.8cm]{GHS05}\ghs[0.8cm]{GHS06}\ghs[0.8cm]{GHS08}}
% ledger: ghs4:K2PtCl4
رباعي كلوروبلاتينات(II) البوتاسيوم: سام إذا ابتُلع، ويسبب تلفًا شديدًا
للعين، ومحسِّس للجهاز التنفسي وللجلد. وتُتداول معقدات البلاتين
بالقفازات تحت ساحبة الأبخرة؛ والسيسبلاتين ونظائره أدوية سامة للخلايا
وتُتداول على هذا الأساس.
\end{safety}

\begin{history}[تاوبه وماركوس]
صنّف هنري تاوبه المعقدات إلى سريعة الاستبدال وخاملة في سنة 1952، وأثبت بتجربته
على الكروم والكوبالت سنة 1953 الآلية عبر الكرة الداخلية؛ 
ونال جائزة نوبل في الكيمياء لسنة 1983. واشتق رودولف ماركوس منذ 1956
نظرية انتقال الإلكترون التي تحمل اسمه، ونال جائزة نوبل
في الكيمياء لسنة 1992.
\end{history}

\section{تمارين}

\begin{exercise}[$\star$]\label{exo:b3:complex-mechanisms:1}
صنّف إلى سريعة الاستبدال أو خاملة، مع التعليل: $\ce{[Cr(H2O)6]^{3+}}$، $\ce{[Co(NH3)6]^{3+}}$، $\ce{[Cr(H2O)6]^{2+}}$،
$\ce{[Cu(H2O)6]^{2+}}$، $\ce{[Zn(H2O)6]^{2+}}$.
\end{exercise}

\begin{exercise}[$\star$]\label{exo:b3:complex-mechanisms:2}
بيّن أن قانون السرعة للآلية D يؤول إلى $v = k_1[\mathrm{ML_5X}]$ عند $[\mathrm Y]$ مرتفع وإلى $v = (k_1k_2/k_{-1})
[\mathrm{ML_5X}][\mathrm Y]/[\mathrm X]$ حين $k_{-1}[\mathrm X] \gg k_2[\mathrm Y]$.
\end{exercise}

\begin{exercise}[$\star$]\label{exo:b3:complex-mechanisms:3}
تنبّأ بإشارتي $\Delta V^{\ddagger}$ و$\Delta^{\ddagger}S^\circ$ لاستبدال من نمط D ولاستبدال من نمط A.
\end{exercise}

\begin{exercise}[$\star$]\label{exo:b3:complex-mechanisms:4}
أعطِ نواتج تفاعل $\ce{[PtCl4]^{2-}}$ مع جزيئي \ce{NH3}، ونواتج تفاعل $\ce{[Pt(NH3)4]^{2+}}$ مع أيوني \ce{Cl-}.
\end{exercise}

\begin{exercise}[$\star\star$]\label{exo:b3:complex-mechanisms:5}
لاستبدال على أيون مائي، $K_{\mathrm{os}} = \qty{2.0}{L.mol^{-1}}$ و$k_i = \qty{3.0e4}{s^{-1}}$ (معطيات
التمرين). احسب $k_{\mathrm{obs}}$ عند $[\mathrm Y] = 0.10$ وعند \qty{1.0}{mol/L}، وحدّه.
\end{exercise}

\begin{exercise}[$\star\star$]\label{exo:b3:complex-mechanisms:6}
يتضاعف ثابت سرعة حين يرتفع الضغط من 0.1 إلى \qty{100}{MPa} عند \qty{298}{K}.
احسب $\Delta V^{\ddagger}$ واستنتج.
\end{exercise}

\begin{exercise}[$\star\star$]\label{exo:b3:complex-mechanisms:7}
في $\ce{[PtCl3(PEt3)]-}$ طول الرابطة Pt--Cl المقابلة للربيطة \ce{PEt3} يساوي \qty{2.42}{\angstrom} وطول الرابطتين الأخريين \qty{2.32}{\angstrom}
(معطيات التمرين). فسّر، وتنبّأ بالكلوريد الذي يُستبدل أولًا.
\end{exercise}

\begin{exercise}[$\star\star$]\label{exo:b3:complex-mechanisms:8}
عدّد الأدلة التجريبية التي تبيّن أن انتقال إلكترون يجري عبر الكرة
الداخلية.
\end{exercise}

\begin{exercise}[$\star\star$]\label{exo:b3:complex-mechanisms:9}
من أجل $\lambda = \qty{80}{kJ/mol}$، احسب $\Delta G^{\ddagger}$ من أجل $\Delta G^\circ = 0$ و$-20$ و$-80$ و\qty{-120}{kJ/mol}.
\end{exercise}

\begin{exercise}[$\star\star\star$]\label{exo:b3:complex-mechanisms:10}
لمزدوجتين ثابتا تبادل ذاتي $k_{11} = 4.0$ و$k_{22} = \qty{2.0e5}{L.mol^{-1}.s^{-1}}$، 
ولتفاعلهما التقاطعي $K_{12} = \num{1.0e4}$ (معطيات التمرين). قدّر $k_{12}$ ($f = 1$).
\end{exercise}

\begin{exercise}[$\star\star\star$]\label{exo:b3:complex-mechanisms:11}
باستعمال نموذج $\sigma$ وحده، احسب طاقات التنشيط لمجال الربائط للأيونات
$d^5$ و$d^7$ و$d^8$ عالية السبين ورتّب \ce{Mn^{2+}} و\ce{Co^{2+}} و\ce{Ni^{2+}} حسب سرعة
تبادل الماء المتوقعة.
\end{exercise}

\begin{exercise}[$\star\star\star$]\label{exo:b3:complex-mechanisms:12}
لماذا يصعب رصد \omterm{def:b3:complex-mechanisms:inverted}{المنطقة المقلوبة} في تفاعلات بين أيونات
تتلاقى بالانتشار في المحلول، وكيف كشفتها الجزيئات الصلبة ذات المانح والمستقبل
المرتبطين؟
\end{exercise}

\section{مسألة: صنع السيسبلاتين لا الترانسبلاتين}

\begin{problem}[{مسألة نهاية الأسبوع --- صنع السيسبلاتين لا الترانسبلاتين:
أيون البلاتين(II) المربع المستوي، وتخليق مخطَّط بمفعول ترانس، 
وتميّه الدواء، ولماذا يُنشَّط داخل الخلايا}]
\label{pb:b3:complex-mechanisms:1}
% ledger: ghs4:K2PtCl4
يُصنع السيسبلاتين، \babelsublr{$cis$-$\ce{[PtCl2(NH3)2]}$}، من $\ce{K2[PtCl4]}$. معطيات المسألة: عند
\qty{37}{\celsius} للتميّه الأول، $\ce{[PtCl2(NH3)2] + H2O -> [PtCl(H2O)(NH3)2]+ + Cl-}$،
ثابت السرعة $k = \qty{9.0e-5}{s^{-1}}$ وثابت التوازن $K = \qty{3.0e-3}{mol/L}$؛ 
وتركيز الكلوريد نحو \qty{100}{mM} في بلازما الدم و\qty{4}{mM} داخل الخلايا.

\medskip\noindent\textbf{الجزء الأول --- البلاتين(II).}
\begin{enumerate}
  \item أعطِ عدد إلكترونات $d$ في \ce{Pt^{2+}}.
  \item لماذا تكون معقداته مربعة مستوية وديامغناطيسية؟
  \item لماذا تكون استبدالاتها ترابطية؟
  \item اكتب قانون السرعة ذا الحدين وفسّر حدّيه.
  \item هل الدواء \omterm{def:b3:complex-mechanisms:labile}{سريع الاستبدال} أم خامل على مقياس زمني قدره يوم؟
\end{enumerate}

\medskip\noindent\textbf{الجزء الثاني --- التخليق.}
\begin{enumerate}[resume]
  \item ماذا يعطي $\ce{K2[PtCl4]}$ مع جزيئي \ce{NH3}؟ ولماذا؟
  \item يعطي هذا الطريق المباشر ناتجًا غير نقي. ويحوّل طريق دارا أولًا
        $\ce{[PtCl4]^{2-}}$ إلى $\ce{[PtI4]^{2-}}$ بفائض من KI. لماذا يفيد اليوديد؟
  \item تعطي إضافة \ce{NH3} المركب \babelsublr{$cis$-$\ce{[PtI2(NH3)2]}$}. فسّر الكيمياء الفراغية.
  \item تُزال اليوديدات بعد ذلك بنترات الفضة في الماء. ماذا يتكوّن؟
  \item تعطي إضافة KCl السيسبلاتين. لماذا لا يتغير التشكيل؟
  \item كيف تصنع الترانسبلاتين بدلًا منه؟
  \item لماذا لا ينفع الترانسبلاتين دواءً، مع أنه يتفاعل هو أيضًا؟
\end{enumerate}

\medskip\noindent\textbf{الجزء الثالث --- التميّه.}
\begin{enumerate}[resume]
  \item احسب عمر النصف للتميّه الأول.
  \item احسب النسبة المتميّهة بعد \qty{2.0}{h} في محلول خالٍ من الكلوريد.
  \item لماذا يكون المعقد المائي هو الصورة النشطة إزاء الحمض النووي؟
  \item هل التميّه تفككي أم ترابطي؟ وكم يكون \omterm{def:b3:complex-mechanisms:activation-volume}{حجم
        تنشيطه}؟
  \item لماذا يُحضَّر محلول الدواء للحقن في محلول ملحي؟
\end{enumerate}

\medskip\noindent\textbf{الجزء الرابع --- البلازما والخلية.}
\begin{enumerate}[resume]
  \item اكتب نسبة المعقد المائي عند التوازن بدلالة
        $[\ce{Cl-}]$.
  \item احسبها في البلازما.
  \item احسبها داخل خلية.
  \item احسب النسبة بينهما.
  \item لماذا يعمل الدواء إذن داخل الخلايا أساسًا؟
  \item ما الربائط الأخرى داخل الخلية التي يمكن أن ترتبط بالبلاتين وتعطّل
        الدواء؟
  \item لماذا تبقى ربيطتا الأمين مرتبطتين طوال الوقت؟
  \item اذكر النتيجة: نسبة كسر المعقد المائي داخل خلية إلى
        كسره في البلازما.
\end{enumerate}
\end{problem}
